本研究提出一項以單張中國傳統國畫圖片為輸入、自動生成 2.5D 動態影片的端到端流程。
由於現有擴散式深度估算模型在國畫風格圖像上容易產生誤判，
本研究引入 Marigold-DC，允許使用者以少量稀疏深度控制點引導估算過程，
並設計專為國畫特性量身打造的複合損失函數 \textit{ChinesePaintingLoss}，
涵蓋筆觸感知平滑損失、拉普拉斯結構損失、留白正則化與輔助排序損失。
取得精煉深度圖後，以 Depth Anything V3 生成的深度為基礎，
透過反投影將其轉為三維點雲，
沿預設攝影機軌跡以光柵化（rasterization）渲染成含視差感的初步點雲影片，
最終輸入 Uni3C 的 PCDController，由視訊擴散模型填補遮擋破洞並維持國畫整體風格，
輸出流暢完整的 2.5D 動態影片。

\begin{flushleft}
{\bf 關鍵字}： 國畫、2.5D 影片、深度估算、擴散模型、點雲渲染
\end{flushleft}
